\section{A different way to pay for missing dimensions}\label{sec:rank}

The source makes the short family span almost the whole target space. The abstract determinant theorem reveals what that hypothesis pays for: completing the family to a basis. It is useful to retain that cost explicitly when designing a new arithmetic application.

\begin{proposition}[Rank and completion cost]\label{prop:rank}
Assume the hypotheses of the upstream corrected-covolume inequality \eqref{eq:covolume-source}. Let the target dimension be $d=n+n'$, with $n>0$, and let $m>0$. Put $C=b_{\rm short}+b_{\rm denom}$. Suppose
\[
 c\le n(A-mR+u),\qquad b_{\rm volume}\le-dA+B.
\]
Then
\begin{equation}\label{eq:rank-cost}
 \boxed{\displaystyle
 mR\le u+\frac{2\delta+B}{n}
             +\frac{d-n}{n}(C-A).}
\end{equation}
No additional sign assumptions on $A,C,R,u,B$ are needed.
\end{proposition}

\begin{proof}
Substitute the two upper bounds into \eqref{eq:covolume-source}:
\[
 0\le -nmR+nu+2\delta+B+(d-n)(C-A).
\]
Rearrange and divide by $n>0$.
\end{proof}

This elementary rearrangement is useful because the deficit is charged by $C-A$, the \emph{excess} completion cost, rather than by $C$ alone. Write $\varrho=n/d$. For an error scale $E_0>0$, suppose
\[
 \varrho\ge\rho>0,\quad
 u\le UE_0,\quad \delta/d\le IE_0,\quad B/d\le JE_0,\quad
 (1-\varrho)(C-A)_+\le KE_0.
\]
Then
\begin{equation}\label{eq:rank-error}
 R\le\frac{U+(2I+J+K)/\rho}{m}\,E_0.
\end{equation}
There are consequently two ways to achieve the needed accuracy: make the missing dimension fraction extremely small, or prove that a completion has excess cost $O(E_0)$. In the latter case a fixed positive fraction of independent target forms can suffice. If $C\le A$, the deficit term is favorable. A uniform lower bound on $n/d$, or a compensating improvement in the other errors, is still needed.

\subsection{An explicit weakening inside the existing proof}

In the source's notation, let
\[
 d=\mathrm{targetDim},\quad A=\mathrm{mainAmount},\quad
 u=fa(\log b)^4,\quad C=2x_1^2b^2+b.
\]
Before using its near-full-rank hypothesis, the proof of \code{integral_forms_long} has already obtained
\begin{equation}\label{eq:source-before-rank}
 0\le c+2\delta+(d-n)C+d(b+u-A).
\end{equation}
The completing vectors come from a globally spanning auxiliary family, with squared-norm exponent $2x_1^2b^2$, scalar denominator-height bound $b$, and the displayed volume-height bound. These arithmetic inputs must be retained; an arbitrary completion is insufficient.\cite{long_source}

Equation~\eqref{eq:source-before-rank} gives
\[
 c+2\delta\ge d(A-b-u)-(d-n)C.
\]
Thus the source's conclusion
\[
                         c+2\delta\ge d(A-3u)
\]
already follows under the completion budget
\begin{equation}\label{eq:weaker-spanning}
                  (d-n)C\le d(2u-b).
\end{equation}
Since the source has $b\le u$, a convenient sufficient condition is $(d-n)C\le du$. The original argument instead obtains the smaller allowance $(d-n)C\le db$ from the loss bound $d-n\le4d/b^2$ and its scale inequalities.

Keeping those same completion and height estimates, it is therefore sufficient to replace that particular loss bound by
\[
 \frac{d-n}{d}\le
      \min\left\{\frac12,\frac{u}{2x_1^2b^2+b}\right\}.
\]
The first term ensures $n\ge d/2$ when the other errors are divided by $n$. With $A\ge0$, $\delta\le dI_0$, $b\le u$, and the short-family upper bound from Proposition~\ref{prop:rank}, this gives
\[
                             mR\le7u+4I_0.
\]
Indeed, use $B=d(b+u)$ in \eqref{eq:rank-cost}, drop the favorable term $-(d-n)A$, and bound the numerator by $d(2I_0+3u)$.

This is a genuine weakening of a sufficient hypothesis in the written source argument, not an improvement to the already proved RH endpoint. For the same expensive completion the missing fraction still tends to zero. For a new target, however, the calculation suggests a distinct research task: construct cheaper integral completing vectors and control their twist scalars, instead of attempting immediately to reproduce the full upstream spanning estimate.

\subsection{What a nonabelian transfer would still need}

The abstract inequality requires injective twist maps, a family and a completion that together form a basis, integral vectors after a common scalar, norm bounds at all complex conjugates, and height bounds at the allowed finite places. The determinant ratio is tied to this same lattice. An abstract frame, a representation-theoretic identity, or a numerical eigenvalue bound does not supply these arithmetic properties.

With the source-like error scale $E_0=fa(\log b)^4$ and $m\asymp\sqrt x$, \eqref{eq:rank-error} has normalized size $O((\log b)^4/\sqrt x)$. Converting this into the final prime remainder still needs the separate comparison of the arithmetic gain with the relevant Frobenius statistic. That is precisely the missing step in the earlier nonabelian criteria.
